\chapter{序論}	% TODO: 章題を記入．題は任意．
\thispagestyle{plain}   % chapterの直後に必ず指定

%TODO: 章の内容を記入．以下はサンプル．

\section{はじめに}
ビデオを使うと\cite{item1}，伝えたい内容を明確に表現できます．[オンライン ビデオ] をクリックすると，追加したいビデオを，それに応じた埋め込みコードの形式で貼り付けできるようになります．キーワードを入力して，文書に最適なビデオをオンラインで検索することもできます．

\section{目的}
ビデオを使うと，伝えたい内容を明確に表現できます．[オンライン ビデオ] をクリックすると，追加したいビデオを，それに応じた埋め込みコードの形式で貼り付けできるようになります．キーワードを入力して，文書に最適なビデオをオンラインで検索することもできます．

\section{章構成}

ビデオを使うと，伝えたい内容を明確に表現できます．[オンライン ビデオ] をクリックすると，追加したいビデオを，それに応じた埋め込みコードの形式で貼り付けできるようになります．キーワードを入力して，文書に最適なビデオをオンラインで検索することもできます．

Word に用意されているヘッダー，フッター，表紙，テキスト ボックス デザインを組み合わせると，プロのようなできばえの文書を作成できます．たとえば，一致する表紙，ヘッダー，サイドバーを追加できます．[挿入] をクリックしてから，それぞれのギャラリーで目的の要素を選んでください．
テーマとスタイルを使って，文書全体の統一感を出すこともできます．[デザイン] をクリックし新しいテーマを選ぶと，図やグラフ，SmartArt グラフィックが新しいテーマに合わせて変わります．スタイルを適用すると，新しいテーマに適合するように見出しが変更されます．

Word では，必要に応じてその場に新しいボタンが表示されるため，効率よく操作を進めることができます．文書内に写真をレイアウトする方法を変更するには，写真をクリックすると，隣にレイアウト オプションのボタンが表示されます．表で作業している場合は，行または列を追加する場所をクリックして，プラス記号をクリックします．

新しい閲覧ビューが導入され，閲覧もさらに便利になりました．文書の一部を折りたたんで，必要な箇所に集中することができます．最後まで読み終わる前に中止する必要がある場合，Word では，たとえ別のデバイスであっても，どこまで読んだかが記憶されます．

\begin{figure}
    \centering
    \includegraphics[width=0.7\textwidth]{fig1.png}
    \caption{図のキャプション}
    \label{fig:my_label}
\end{figure}

ビデオを使うと，伝えたい内容を明確に表現できます．[オンライン ビデオ] をクリックすると，追加したいビデオを，それに応じた埋め込みコードの形式で貼り付けできるようになります．キーワードを入力して，文書に最適なビデオをオンラインで検索することもできます．

Word に用意されているヘッダー，フッター，表紙，テキスト ボックス デザインを組み合わせると，プロのようなできばえの文書を作成できます．たとえば，一致する表紙，ヘッダー，サイドバーを追加できます．[挿入] をクリックしてから，それぞれのギャラリーで目的の要素を選んでください．
テーマとスタイルを使って，文書全体の統一感を出すこともできます．[デザイン] をクリックし新しいテーマを選ぶと，図やグラフ，SmartArt グラフィックが新しいテーマに合わせて変わります．スタイルを適用すると，新しいテーマに適合するように見出しが変更されます．

Word では，必要に応じてその場に新しいボタンが表示されるため，効率よく操作を進めることができます．文書内に写真をレイアウトする方法を変更するには，写真をクリックすると，隣にレイアウト オプションのボタンが表示されます．表で作業している場合は，行または列を追加する場所をクリックして，プラス記号をクリックします．

新しい閲覧ビューが導入され，閲覧もさらに便利になりました．文書の一部を折りたたんで，必要な箇所に集中することができます．最後まで読み終わる前に中止する必要がある場合，Word では，たとえ別のデバイスであっても，どこまで読んだかが記憶されます．

\begin{table}[htb]
    \caption{表のキャプション}
    \centering
    \begin{tabular}{lcr}
        \hline
        head1 & head2 & head3 \\ 
        \hline\hline
        123 & 456 & 789\\
        \hline
    \end{tabular}
    \label{tab:my_label}
\end{table}

ビデオを使うと，伝えたい内容を明確に表現できます．[オンライン ビデオ] をクリックすると，追加したいビデオを，それに応じた埋め込みコードの形式で貼り付けできるようになります．キーワードを入力して，文書に最適なビデオをオンラインで検索することもできます．

Word に用意されているヘッダー，フッター，表紙，テキスト ボックス デザインを組み合わせると，プロのようなできばえの文書を作成できます．たとえば，一致する表紙，ヘッダー，サイドバーを追加できます．[挿入] をクリックしてから，それぞれのギャラリーで目的の要素を選んでください．
テーマとスタイルを使って，文書全体の統一感を出すこともできます．[デザイン] をクリックし新しいテーマを選ぶと，図やグラフ，SmartArt グラフィックが新しいテーマに合わせて変わります．スタイルを適用すると，新しいテーマに適合するように見出しが変更されます．

Word では，必要に応じてその場に新しいボタンが表示されるため，効率よく操作を進めることができます．文書内に写真をレイアウトする方法を変更するには，写真をクリックすると，隣にレイアウト オプションのボタンが表示されます．表で作業している場合は，行または列を追加する場所をクリックして，プラス記号をクリックします．
新しい閲覧ビューが導入され，閲覧もさらに便利になりました．文書の一部を折りたたんで，必要な箇所に集中することができます．最後まで読み終わる前に中止する必要がある場合，Word では，たとえ別のデバイスであっても，どこまで読んだかが記憶されます．

ビデオを使うと，伝えたい内容を明確に表現できます．[オンライン ビデオ] をクリックすると，追加したいビデオを，それに応じた埋め込みコードの形式で貼り付けできるようになります．キーワードを入力して，文書に最適なビデオをオンラインで検索することもできます．

Word に用意されているヘッダー，フッター，表紙，テキスト ボックス デザインを組み合わせると，プロのようなできばえの文書を作成できます．たとえば，一致する表紙，ヘッダー，サイドバーを追加できます．[挿入] をクリックしてから，それぞれのギャラリーで目的の要素を選んでください．
テーマとスタイルを使って，文書全体の統一感を出すこともできます．[デザイン] をクリックし新しいテーマを選ぶと，図やグラフ，SmartArt グラフィックが新しいテーマに合わせて変わります．スタイルを適用すると，新しいテーマに適合するように見出しが変更されます．

Word では，必要に応じてその場に新しいボタンが表示されるため，効率よく操作を進めることができます．文書内に写真をレイアウトする方法を変更するには，写真をクリックすると，隣にレイアウト オプションのボタンが表示されます．表で作業している場合は，行または列を追加する場所をクリックして，プラス記号をクリックします．

新しい閲覧ビューが導入され，閲覧もさらに便利になりました．文書の一部を折りたたんで，必要な箇所に集中することができます．最後まで読み終わる前に中止する必要がある場合，Word では，たとえ別のデバイスであっても，どこまで読んだかが記憶されます．
